\Needspace{15\baselineskip}\Article{9}

Als Zweites betrachten wir die folgende Reihe:

\[
\begin{aligned}
&1+x^{\frac12}\frac{1-x^m}{1-x}+x\frac{(1-x^m)(1-x^{m-1})}{(1-x)(1-xx)}\\
&\qquad+x^{\frac32}\frac{(1-x^m)(1-x^{m-1})(1-x^{m-2})}{(1-x)(1-xx)(1-x^3)}+\text{usw.}
\end{aligned}
\]

oder

\[
\begin{aligned}
1+x^{\frac12}(m,1)+x(m,2)+x^{\frac32}(m,3)+xx(m,4)+\text{usw.}
\end{aligned}
\]

die wir mit $F(x,m)$ bezeichnen.\quad Wir beschränken diese Untersuchung auf den Fall, dass $m$ eine positive ganze Zahl ist, sodass auch diese Reihe stets
\SourceLeaf{31}
mit ihrem Glied an der Stelle $m+1$ abbricht, wobei $=x^{\frac12m}(m,m)$ gilt.\quad Da

\[
\begin{aligned}
(m,m)=1,\quad(m,m-1)=(m,1),\quad(m,m-2)=(m,2),\;\text{usw.}
\end{aligned}
\]

ist, lässt sich die Reihe auch so darstellen:

\[
\begin{aligned}
F(x,m)&=x^{\frac12m}+x^{\frac12(m-1)}(m,1)\\&\qquad+x^{\frac12(m-2)}(m,2)+x^{\frac12(m-3)}(m,3)+\text{usw.}
\end{aligned}
\]

Daraus ergibt sich

\[
\begin{aligned}
(1+x^{\frac12m+\frac12})F(x,m)&=1+x^{\frac12}(m,1)+x(m,2)\\&\qquad+x^{\frac32}(m,3)+\text{usw.}\\
&+x^{\frac12}\pdot x^m+x\pdot x^{m-1}(m,1)\\&\qquad+x^{\frac32}\pdot x^{m-2}(m,2)+\text{usw.}
\end{aligned}
\]

Da nun (Art. 5, II)

\[
\begin{aligned}
(m,1)+x^m&=(m+1,1)\\
(m,2)+x^{m-1}(m,1)&=(m+1,2)\\
(m,3)+x^{m-2}(m,2)&=(m+1,3),\;\text{usw.}
\end{aligned}
\]

gilt, erhalten wir

\[
\begin{aligned}
(1+x^{\frac12m+\frac12})F(x,m)=F(x,m+1)\ .\; .\; .\; .\; .\; .\; .\ [3]
\end{aligned}
\]

Es ist aber $F(x,0)=1$; daher gilt

\[
\begin{aligned}
F(x,1)&=1+x^{\frac12}\\
F(x,2)&=(1+x^{\frac12})(1+x)\\
F(x,3)&=(1+x^{\frac12})(1+x)(1+x^{\frac32}),\;\text{usw.}
\end{aligned}
\]

oder allgemein

\[
\begin{aligned}
F(x,m)&=(1+x^{\frac12})(1+x)(1+x^{\frac32})\\&\qquad\dotsfour(1+x^{\frac12m})\ .\; .\; .\; .\quad[4]
\end{aligned}
\]

\Needspace{12\baselineskip}\Article{10}

Nach diesen vorbereitenden Untersuchungen wollen wir unserem Ziel nun näherkommen.\quad Für einen Primzahlwert von $n$ sind die Quadrate $1$, $4$, $9\dotsfour(\frac12(n-1))^2$ untereinander sämtlich inkongruent modulo $n$. Ihre kleinsten Reste nach diesem Modul müssen daher mit den Zahlen $a$ übereinstimmen, und somit ist

\[
\begin{aligned}
\Sigma\cos ak\omega&=\cos k\omega+\cos4k\omega+\cos9k\omega\\&\qquad+\text{usw.}+\cos(\tfrac12(n-1))^2k\omega\\
\Sigma\sin ak\omega&=\sin k\omega+\sin4k\omega+\sin9k\omega\\&\qquad+\text{usw.}+\sin(\tfrac12(n-1))^2k\omega
\end{aligned}
\]

Da ferner dieselben Quadrate $1$, $4$, $9\dotsfour(\frac12(n-1))^2$ in umgekehrter Reihenfolge zu $(\frac12(n+1))^2$, $(\frac12(n+3))^2$, $(\frac12(n+5))^2\dotsfive(n-1)^2$ kongruent sind, gilt auch

\SourceLeaf{32}

\[
\begin{aligned}
\Sigma\cos ak\omega&=\cos(\tfrac12(n+1))^2k\omega\\&\qquad+\cos(\tfrac12(n+3))^2k\omega+\text{usw.}+\cos(n-1)^2k\omega\\
\Sigma\sin ak\omega&=\sin(\tfrac12(n+1))^2k\omega\\&\qquad+\sin(\tfrac12(n+3))^2k\omega+\text{usw.}+\sin(n-1)^2k\omega
\end{aligned}
\]

Wir setzen also

\[
\begin{aligned}
T&=1+\cos k\omega+\cos4k\omega+\cos9k\omega\\&\qquad+\text{usw.}+\cos(n-1)^2k\omega\\
U&=\phantom{1+}\sin k\omega+\sin4k\omega+\sin9k\omega\\&\qquad+\text{usw.}+\sin(n-1)^2k\omega
\end{aligned}
\]

Dann ist

\[
\begin{aligned}
1+2\Sigma\cos ak\omega&=T\\
2\Sigma\sin ak\omega&=U
\end{aligned}
\]

Daraus ist ersichtlich, dass die in Art. 1 angegebenen Summationen von der Summation der Reihen $T$ und $U$ abhängen. Wir lassen daher die ersteren beiseite, richten unsere Untersuchung auf diese Reihen und führen sie so allgemein durch, dass sie sowohl Primzahlwerte von $n$ als auch beliebige zusammengesetzte Werte umfasst.\quad Wir setzen dabei voraus, dass die Zahl $k$ zu $n$ teilerfremd ist: Der Fall, dass $k$ und $n$ einen gemeinsamen Teiler haben, lässt sich ohne Mühe auf diesen zurückführen.

\Needspace{12\baselineskip}\Article{11}

Wir bezeichnen die imaginäre Größe $\surd{-1}$ mit $i$ und setzen

\[
\begin{aligned}
\cos k\omega+i\sin k\omega=r
\end{aligned}
\]

Dann ist $r^n=1$, oder $r$ ist eine Wurzel der Gleichung $r^n-1=0$.\quad Man sieht leicht, dass sämtliche Zahlen $k$, $2k$, $3k\dotsfour(n-1)k$ nicht durch $n$ teilbar und untereinander inkongruent modulo $n$ sind. Daher sind die Potenzen von $r$

\[
\begin{aligned}
1,\;r,\;rr,\;r^3\dotsfour r^{n-1}
\end{aligned}
\]

sämtlich verschieden, während jede von ihnen auch der Gleichung $x^n-1=0$ genügt. Aus diesem Grund stellen diese Potenzen \emph{alle} Wurzeln der Gleichung $x^n-1=0$ dar.

Diese Schlüsse würden nicht gelten, wenn $k$ mit $n$ einen gemeinsamen Teiler hätte. Wäre nämlich $\nu$ ein solcher gemeinsamer Teiler, so wäre $k\pdot\dfrac n\nu$ durch $n$ teilbar und damit eine niedrigere Potenz als $r^n$, nämlich $r^{\frac n\nu}$, gleich der Einheit.\quad In diesem Fall stellen die Potenzen von $r$ also höchstens $\dfrac n\nu$ Wurzeln der Gleichung $x^n-1=0$ dar, und tatsächlich ergeben sie ebenso viele verschiedene Wurzeln, wenn $\nu$ der \emph{größte} gemeinsame Teiler der Zahlen\SourceLeaf{33}
$k$, $n$ ist.\quad In unserem Fall, in dem $k$ und $n$ als teilerfremd vorausgesetzt sind, lässt sich $r$ zweckmäßig eine \emph{eigentliche Wurzel} der Gleichung $x^n-1=0$ nennen. \mbox{Im anderen Fall} dagegen, in dem $k$ und $n$ einen (größten) gemeinsamen Teiler $\nu$ hätten, würde $r$ eine \emph{uneigentliche Wurzel} jener Gleichung heißen. Dann wäre aber dieselbe $r$ offenbar eine eigentliche Wurzel der Gleichung $x^{\frac n\nu}-1=0$.\quad Die einfachste uneigentliche Wurzel ist die Einheit, und im Fall, dass $n$ eine Primzahl ist, gibt es überhaupt keine anderen uneigentlichen Wurzeln.

\Needspace{13\baselineskip}\Article{12}

Setzen wir nun

\[
\begin{aligned}
W=1+r+r^4+r^9+\text{usw.}+\text{usw.}+r^{(n-1)^2}
\end{aligned}
\]

so ist offenbar $W=T+iU$; somit ist $T$ der reelle Teil von $W$, und $U$ geht aus dem imaginären Teil von $W$ hervor, wenn der Faktor $i$ weggelassen wird.\quad Damit führt sich die ganze Frage auf die Bestimmung der Summe $W$ zurück. Hierzu lässt sich entweder die in Art. 6 betrachtete Reihe verwenden oder diejenige, deren Summation wir in Art. 9 gezeigt haben; die erstere ist jedoch weniger geeignet, wenn $n$ gerade ist.\quad Dennoch hoffen wir, dass es den Lesern willkommen sein wird, wenn wir den Fall, dass $n$ ungerade ist, nach zwei Methoden behandeln.

Zunächst sei also $n$ eine ungerade Zahl, $r$ bezeichne eine beliebige eigentliche Wurzel der Gleichung $x^n-1=0$, und in der Funktion $f(x,m)$ werde $x=r$ und $m=n-1$ gesetzt.\quad Daraus ergibt sich offenbar

\[
\begin{aligned}
\frac{1-x^m}{1-x}&=\frac{1-r^{-1}}{1-r}=-r^{-1}\\
\frac{1-x^{m-1}}{1-xx}&=\frac{1-r^{-2}}{1-rr}=-r^{-2}\\
\frac{1-x^{m-2}}{1-x^3}&=\frac{1-r^{-3}}{1-r^3}=-r^{-3},\;\text{usw.}
\end{aligned}
\]

bis hin zu

\[
\begin{aligned}
\frac{1-x}{1-x^m}=\frac{1-r^{-m}}{1-r^m}=-r^{-m}
\end{aligned}
\]

(Es ist nicht überflüssig, darauf hinzuweisen, dass diese Gleichungen nur insoweit gelten, als $r$ als eigentliche Wurzel vorausgesetzt ist: Wäre $r$ eine uneigentliche Wurzel, so würden bei einigen dieser Brüche Zähler und Nenner zugleich verschwinden, und die Brüche würden somit unbestimmt).

Daraus leiten wir die folgende Gleichung ab:

\SourceLeaf{34}

\[
\begin{aligned}
f(r,n-1)&=1+r^{-1}+r^{-3}+r^{-6}+\text{usw.}+r^{-\frac12(n-1)n}\\
&=(1-r)(1-r^3)(1-r^5)\dotsfive(1-r^{n-2})
\end{aligned}
\]

\Needspace{14\baselineskip}Dieselbe Gleichung gilt weiterhin, wenn wir $r$ durch $r^\lambda$ ersetzen, wobei $\lambda$ eine beliebige zu $n$ teilerfremde ganze Zahl bezeichnet: Dann ist auch $r^\lambda$ eine eigentliche Wurzel der Gleichung $x^n-1=0$.\quad Ersetzen wir also $r$ durch $r^{n-2}$ oder, was dasselbe ist, durch $r^{-2}$, so erhalten wir

\[
\begin{aligned}
&1+r^2+r^6+r^{12}+\text{usw.}+r^{(n-1)n}\\&\qquad=(1-r^{-2})(1-r^{-6})(1-r^{-10})\ldots(1-r^{-2(n-2)})
\end{aligned}
\]

Multiplizieren wir beide Seiten dieser Gleichung mit

\[
\begin{aligned}
r\pdot r^3\pdot r^5\dotsfour r^{(n-2)}=r^{\frac14(n-1)^2}
\end{aligned}
\]

\Needspace{13\baselineskip}so ergibt sich wegen

\[
\begin{aligned}
r^{2+\frac14(n-1)^2}&=r^{\frac14(n-3)^2},\\r^{(n-1)n+\frac14\,n-1)^2}=r^{\frac14(n+1)^2}\\
r^{6+\frac14(n-1)^2}&=r^{\frac14(n-5)^2},\\r^{(n-2)(n-1)+\frac14(n-1)^2}=r^{\frac14(n+3)^2}\\
r^{12+\frac14(n-1)^2}&=r^{\frac14(n-7)^2},\\r^{(n-3)(n-2)+\frac14(n-1)^2}=r^{\frac14(n+5)^2},\;\text{usw.}
\end{aligned}
\]

die folgende Gleichung:

\[
\begin{aligned}
&r^{\frac14(n-1)^2}+r^{\frac14(n-3)^2}+r^{\frac14(n-5)^2}+\text{usw.}+r+1\\
&+r^{\frac14(n+1)^2}+r^{\frac14(n+3)^2}+r^{\frac14(n+5)^2}\\&\qquad+\text{usw.}+r^{\frac14(2n-2)^2}\\
&\qquad=(r-r^{-1})(r^3-r^{-3})(r^5-r^{-5})\dotsfour(r^{n-2}-r^{-n+2})
\end{aligned}
\]

\Needspace{7\baselineskip}oder, wenn wir die Glieder der linken Seite anders anordnen,

\[
\begin{aligned}
&1+r+r^4+\text{usw.}+r^{(n-1)^2}\\&\qquad=(r-r^{-1})(r^3-r^{-3})\ldots(r^{n-2}-r^{-n+2})\quad\dotstwo\quad[5]
\end{aligned}
\]
